\chapter{Zusammenfassung und Ausblick}
\label{ch:Zusammenfassung}

Die aufbereiteten Sensordaten wurden mit vier getrennten Verfahren untersucht:
Ein LSTM und ein kompakter Transformer bewerten Abweichungen von ihrer
Prognose, ein Autoencoder bewertet Rekonstruktionsfehler und eine robuste
parameterfreie Baseline die Abweichung vom gemessenen Wert 24 Stunden zuvor.
Die drei neuronalen Modelle lernen zusätzlich die Verweildauer; sämtliche
Alarme und künstlichen Testanomalien beziehen sich aber ausschließlich auf die
Passantenfrequenz. Die Verfahren verwenden dieselben Zeitabschnitte und
Entscheidungskriterien, ohne ihre Anomaliewerte miteinander zu vermischen.

Der LSTM erreicht auf der gemeinsamen operativen Testpopulation den höchsten
Punkt-F1, die höchste AP und die beste Rauschrobustheit. Der Transformer
prognostiziert genauer und erreicht den höchsten Ereignis-Recall sowie knapp
den höchsten Ereignis-F1, reagiert aber stärker auf Rauschen. Die saisonale
Baseline benötigt keine trainierbaren Parameter, bleibt nach Einbeziehung
ihrer nicht bewertbaren Stunden aber hinter beiden Prognosemodellen. Der
Autoencoder ist kleiner, bleibt bei den Gesamtkennzahlen jedoch ebenfalls
dahinter. Die zusätzliche Attention führt somit zu einem konkreten
Ereignisvorteil, aber nicht zu einer allgemeinen Überlegenheit.
Als nächste Schritte sind vor allem die vollständige manuelle Prüfung der nach
dem letzten Training aktuellen Kandidaten, die Klärung der geplanten
Sensorlaufzeiten, längere Messreihen und ein vollständig neuer Testzeitraum
sinnvoll. Erst dann sind echte Rolling-Origin-Auswertungen und eine erneute
Abstimmung von Fensterlänge und Transformer-Größe belastbar. Insbesondere der
kleine Vorsprung bei den Ereignismetriken muss auf längeren Zeiträumen
bestätigt werden.

